% Generated draft: requires mathematical and source review.
\documentclass[11pt]{article}
\usepackage[margin=25mm]{geometry}
\usepackage{fontspec}
\setmainfont{Latin Modern Roman}
\usepackage{amsmath,amssymb,mathtools,unicode-math}
\setmathfont{Latin Modern Math}
\usepackage{enumitem,needspace,xurl,hyperref}
\setlength{\parindent}{0pt}
\setlength{\parskip}{5pt}
\setlength{\emergencystretch}{3em}
\newcommand{\N}{\mathbb{N}}
\newcommand{\Z}{\mathbb{Z}}
\newcommand{\Q}{\mathbb{Q}}
\newcommand{\R}{\mathbb{R}}
\newcommand{\C}{\mathbb{C}}
\newcommand{\F}{\mathbb{F}}
\newcommand{\A}{\mathbb{A}}
\newcommand{\PP}{\mathbb{P}}
\newcommand{\E}{\mathbb{E}}
\newcommand{\eps}{\varepsilon}
\newcommand{\dd}{\,\mathrm{d}}
\begin{document}
\section*{UnsolvedMath ID: 10400181}
\textbf{Source problem code:} catalogue entry
\subsection{Definitions and mathematical statement}
Let \(L\subset S^3\) be a framed link, and let \(M_L\) be the three-manifold obtained by surgery on \(L\), using the chosen orientation and surgery conventions. For each finite-type invariant \(v\) of three-manifolds in the source's intended category, define the induced framed-link invariant
\[f_v(L)=v(M_L)\]
whenever \(M_L\) lies in its domain. Polyak asks for a Gauss-diagram formula for \(f_v\), for \emph{each} such \(v\).

A Gauss diagram of an oriented link diagram consists of one parametrizing circle per component, with a signed arrow for each crossing, directed according to a fixed over/under convention; framing information must additionally be recorded or encoded by the diagram convention. An explicit formula must give a rule evaluating this data to \(v(M_L)\), and must be independent of presentation choices, including any auxiliary component orientations and moves preserving the surgered manifold. In a usual arrow-counting formulation, the rule would specify diagrams \(D\), coefficients and a pairing \(\langle D,G(L)\rangle\) counting signed embeddings of their arrows into the Gauss diagram, together with the required treatment of framings and surgery data. The excerpt does not fix a formula class, coefficient ring or allowed operations, so it does not justify imposing a particular finite linear combination, a degree bound, or prohibiting linking-matrix terms.

The meaning of finite type for three-manifold invariants, their manifold domain and permitted surgery links must be supplied; these are not determined by ordinary finite type of link invariants. The request is uniform over the chosen class of finite-type invariants, not merely a formula for a selected low-degree example.
\subsection{Short English statement}
Express every finite-type three-manifold invariant as a formula on Gauss diagrams of framed surgery links. The finite-type theory, admissible surgeries and precise formula format are not specified in the excerpt.
\subsection{Sources}
\begin{itemize}
\item Local UnsolvedMath registry, permanent ID \(10400181\), exact source code \texttt{catalogue entry}.
\item M. Polyak, attribution in Ohtsuki, \emph{Problems on invariants of knots and 3-manifolds} (2002), Problem 10.6, PDF page 145: \url{https://msp.org/gtm/2002/04/gtm-2002-04-024s.pdf}. Registry access 2026-07-29, extraction pdf-font-structure.
\end{itemize}
\end{document}
